% Requires in preamble: \usepackage{booktabs, longtable, pdflscape, amsmath, cleveref}
\subsection{Size \& Composition}\label{sec:size_compo}
According to Statbel's 2025 cadastral building-stock dataset, Belgium comprised
4{,}668{,}425 buildings and 5{,}827{,}823 dwellings on 1~January 2025
\cite{statbel_building_stock_2025}. As these figures derive from the cadastre,
they cover the entire Belgian building stock. The stock is regionally
concentrated in Flanders, which holds 58.5\% of dwellings, followed by Wallonia
(31.1\%) and the Brussels-Capital Region (10.5\%). Grouping dwellings by the type
of building in which they are located, 65.7\% are in house-type buildings, 29.3\%
in apartment buildings, and the remaining 5.0\% in commercial or other building
categories. The Belgian residential stock is therefore predominantly
house-based; apartments nonetheless account for close to one-third of all
dwellings and constitute the majority of the Brussels stock.

\subsubsection*{Dwelling types}
Because Statbel reports dwellings by building type, its categories can be mapped
onto the dwelling types of the Belgian TABULA/VITO housing typology
\cite{tabula_be_vito}, as summarised in \Cref{tab:dwelling_type_mapping}. The
apartment stock (Statbel category R4) is divided equally between the enclosed and
exposed apartment archetypes, since Statbel does not record an apartment's
position within its building.

\begin{table}[ht]
  \centering
  \caption{Mapping of Statbel building categories onto the five Belgian
           TABULA/VITO dwelling types \cite{tabula_be_vito,statbel_building_stock_2025}.}
  \label{tab:dwelling_type_mapping}
  \begin{tabular}{lll}
    \toprule
    TABULA/VITO dwelling type & Statbel building category & Code \\
    \midrule
    Single-family detached      & Detached/open house, farm, castle & R3 \\
    Single-family semi-detached & Semi-detached house               & R2 \\
    Single-family terraced      & Terraced / closed house           & R1 \\
    Enclosed apartment          & Apartment building (half of R4)   & R4 \\
    Exposed apartment           & Apartment building (half of R4)   & R4 \\
    \bottomrule
  \end{tabular}
\end{table}

These five dwelling-type archetypes provide a first-order approximation of the
Belgian dwelling stock.

\subsubsection*{Construction period}
Belgium's occupied residential stock is old. In 2021, 22\% of occupied
conventional dwellings were located in buildings completed before 1919, whereas
only about 8\% were in buildings completed from 2011 onwards
\cite{statbel_census_construction_2021}. A pronounced regional contrast is also
present: Wallonia has by far the largest pre-1919 share (38\%), while Flanders
holds the highest shares of dwellings built after 1991. Brussels displays a
particularly old profile, with only about 13\% of dwellings built from 1991
onwards and 5\% from 2011 onwards \cite{statbel_census_construction_2021}.

The Belgian TABULA/VITO housing typology is used to aggregate these census bands
into five construction-period classes that track successive shifts in envelope
and insulation practice, as interpreted in
\Cref{tab:construction_period_classes}.

\begin{table}[ht]
  \centering
  \caption{TABULA/VITO construction-period classes and their thermal
           interpretation \cite{tabula_be_vito}.}
  \label{tab:construction_period_classes}
  \begin{tabular}{ll}
    \toprule
    TABULA age class & Interpretation \\
    \midrule
    Before 1946 & Pre-war stock; weak original envelope \\
    1946--1970  & Post-war stock; mostly before strong insulation practice \\
    1971--1990  & Oil-crisis / early insulation transition \\
    1991--2005  & First building-energy regulation period \\
    After 2005  & Post-EPBD / stronger energy-performance period \\
    \bottomrule
  \end{tabular}
\end{table}

Crossing the five age classes in \Cref{tab:construction_period_classes} with the
five dwelling types in \Cref{tab:dwelling_type_mapping} yields the 25 building
archetypes that define the Belgian dwelling stock in this model.

\subsection{Physical Building Parameters}
A dwelling's physical parameters comprise principally its construction elements,
its geometry, its insulation (U-values and infiltration), and its
heating, domestic-hot-water and ventilation systems.

The Belgian TABULA/VITO national housing typology is used to parametrise the 25
archetypes defined in \Cref{sec:size_compo} \cite{tabula_be_vito}. Geometry is
taken from Annex~2, Table~19, which directly lists the 25 national dwelling-unit
types; envelope U-values are assigned by construction period from Table~10; and
airtightness is assigned as $v_{50}$ by dwelling type and age class from
Table~9. This procedure keeps stock weighting separate from physical-parameter
assignment: Statbel determines how frequent each archetype is, while TABULA/VITO
supplies the representative geometry, envelope quality and airtightness. The same
age-class construction-element U-values are applied across dwelling types,
following the TABULA/VITO construction-element sub-typology; differences between
detached, semi-detached, terraced and apartment dwellings are instead captured
through the TABULA geometry, i.e.\ through different exposed envelope areas.
Combining these inputs yields the base archetype matrix reported in
\Cref{tab:base_physical_archetype_matrix}.

\input{tables/base_physical_archetype_matrix.tex}

\subsubsection*{Regional stock weight}
The cadastral data are used to assign a stock share to each archetype and then to
split it by region. Each regional archetype share is obtained as the product of
the regional dwelling-type share and the regional construction-period share,
under a type--period independence assumption; the modelled types R1--R4 are
renormalised within the regional stock and the residual R5/R6 dwellings are
excluded from the regional denominator. A 50/50 split is applied to the two
apartment types, as Statbel's R4 category does not identify an apartment's
position within its building (enclosed versus exposed); the equal division is
therefore an explicit balancing assumption rather than an observed distribution.
